% titel: Blatt 2, A2: Molare Masse von Verbindungen
% skript: Mortimer Kap. 4.1 (Molmasse, Stoffmenge)

% Teil I: Werkzeuge, kein Lösungsweg.

Die molare Masse $M$ einer Verbindung ergibt sich aus der Summe der molaren
Massen aller Atome in der Summenformel. Folgende Punkte sind beim Rechnen zu
beachten.

\begin{werkzeugkasten}[Molare Masse einer Verbindung]
Für jede Atomsorte in der Formel zaehlt man die Anzahl $n$ der Atome und
multipliziert mit der molaren Masse des Elements:
\begin{equation}
  M(\mathrm{Verbindung})
  = \sum_{\text{Atomsorten}} n_{\text{Sorte}} \cdot M_{\text{Atom}}.
\end{equation}
Beispielhaftes Aufaddieren: $\mathrm{MgSO_4 \cdot 7\,H_2O}$ enthaelt pro
Formeleinheit je ein $\mathrm{Mg}$, $\mathrm{S}$, $4\cdot\mathrm{O}$ im Sulfat
plus $7$ Wassermoleküle ($7\cdot 2\,\mathrm{H}$, $7\cdot 1\,\mathrm{O}$).
Eingefangenenfalls Kristallwasser mitzählen --- es gehört zur Masse.

Typische Werte (gerundet): $\mathrm{H} \approx \qty{1.0}{g/mol}$,
$\mathrm{C} \approx \qty{12.0}{g/mol}$, $\mathrm{O} \approx \qty{16.0}{g/mol}$,
$\mathrm{Na} \approx \qty{23.0}{g/mol}$, $\mathrm{P} \approx \qty{31.0}{g/mol}$,
$\mathrm{S} \approx \qty{32.1}{g/mol}$, $\mathrm{Cl} \approx \qty{35.5}{g/mol}$,
$\mathrm{K} \approx \qty{39.1}{g/mol}$, $\mathrm{Cr} \approx \qty{52.0}{g/mol}$,
$\mathrm{Mn} \approx \qty{54.9}{g/mol}$, $\mathrm{Mg} \approx \qty{24.3}{g/mol}$.
\textit{Skript:} Mortimer Kap. 4.1 (Stoffmenge, Molmasse), Übung 2/2.
\end{werkzeugkasten}

\begin{achtung}[Darauf achten]
\begin{itemize}
  \item Die Anzahl jedes Elements aus der Formel sauber aufschreiben, bevor man
        multipliziert --- besonders bei Klammern wie
        $\mathrm{Cr_2(SO_4)_3}$ (da stecken $2$ Cr, $3$ S, $12$ O drin).
  \item Hydrate: $7\,H_2O$ vollstaendig in die Masse einrechnen.
  \item Moleküle wie $\mathrm{C_4H_{11}N}$ einfach nach Elementsymbolen
        abzaehlen.
  \item Einheit durchgehend $\mathrm{g/mol}$; nie den Faktor vergessen, dass
        die Menge pro $1$ mol gerechnet wird.
\end{itemize}
\end{achtung}
